\documentclass[11pt]{article}
\usepackage[margin=0.9in]{geometry}
\usepackage{amsmath,amssymb,amsthm,booktabs,mathrsfs}
\usepackage[T1]{fontenc}
\usepackage{lmodern,microtype}
\usepackage[colorlinks=true,linkcolor=blue,urlcolor=blue]{hyperref}
\newtheorem{theorem}{Theorem}[section]
\newtheorem{lemma}[theorem]{Lemma}
\newtheorem{proposition}[theorem]{Proposition}
\newtheorem{corollary}[theorem]{Corollary}
\theoremstyle{remark}
\newtheorem{remark}[theorem]{Remark}
\newcommand{\norm}[1]{\left\|#1\right\|}
\title{Toward a four-dimensional Yang--Mills construction\\
\large Notebook 29: total Ward extraction and classical spatial gain}
\author{Research working document}
\date{17 October 2026}
\begin{document}
\maketitle
\begin{abstract}
We continue the existence construction by treating the low field degrees
left open by the high-degree estimate of Notebook 28. The actual local
macrostep preserves output translations and the hypercubic group, under
the selected averaging map's stated permutation axiom. Its total
correction obeys the nonlinear gauge identities. Uniform kernel
locality permits a joint subsequential limit of vacuum-normalized
analytic interactions and coherent classical background families. There, the
quadratic, cubic and quartic marginal vertices are all fixed by one
coefficient multiplying a normalized gauge-invariant reference action.
The actual homogeneous classical action has the required normalization,
so this extraction has an exact algebraic coupling rebase.

We also propagate complete first-gradient H\"older regularity through
the ordinary averaging prefix, including strict support families.
Using total moment cancellations before assigning carriers, an exact
lattice factorization puts one difference on each quadratic factor.
Together with cubic and quartic spatial subtraction, this gives a
connected-activity bound of order \(C(J)L^{-\alpha}\) for the canceled
low-degree background, under the specified coherent-family and
infinite-kernel hypotheses. The coefficient \(C(J)\) is retained:
it cannot be defeated by a ratio chosen before \(J\).
Two ordinary minimum comparisons prove a classical background
semigroup modulo gauge. This allows a finite endpoint ratio
\(L=R^q\) to be chosen after fixing the elementary ratio \(R\)
and grouping \(J\).
An activity space modulo equality of normalized totals is Banach.
The full-reference extraction is bounded independently of the
small analytic radius, and classical background substitution
descends to this quotient. Combining the low- and high-degree
estimates gives a classical contraction under the stated joint
map and support hypotheses.

These results give a gauge-compatible marginal classification,
a quantitative classical localization of its remainder, and a
complete space in which the representation can change. They do
not yet establish that the actual quantum update descends to
this space with the required finite-regulator compatibility,
the required coupling trajectory,
or the remaining large-field and continuum construction.
\end{abstract}

The objective remains a complete existence proof for a nontrivial
four-dimensional Yang--Mills theory. The preceding notebooks establish
the homogeneous classical branch and a local quantum macrostep.
This notebook distinguishes exact finite-regulator local identities
from jointly extracted infinite analytic interactions and from the
operator estimates that require those coherent limiting families.
The compactness argument connects these levels at fixed parameters;
it does not provide a universal quantum update on all such limits.
Gauge invariance is required of the total functional,
not of each chosen activity coefficient.

